\subsection{Canonical support and class VII normalization models}

\begin{statement}{C.2.UN.17}{Proved}{Canonical support forces the numerically anticanonical class VII case}
Let $f:U\to\Delta$ be a proper holomorphic map from a smooth
complex threefold, with connected fibres and smooth two-torus
fibres over $\Delta^*$. Write its full central fibre as
$F=\sum_i m_iS_i$. For each component let
$h_i:Z_i\to S_i$ be its normalization followed by minimal
resolution, and let $D_i\ge0$ be the intrinsic resolution
conductor of A.5.3. There are nonnegative integers $a_j$ with
$K_U\simeq\OO_U(\sum_j a_jS_j)$. Put
$E_{ij}=h_i^*(S_j)$ for $j\ne i$, using the effective
Cartier pullback. The exact line-bundle identity is
\[
 K_{Z_i}^{\otimes m_i}\otimes\OO_{Z_i}(m_iD_i)
 \simeq
 \OO_{Z_i}\!\left(
  \sum_{j\ne i}\bigl[m_i a_j-m_j(a_i+1)\bigr]E_{ij}
 \right).                                                \tag{C.CA.1}
\]
In particular a positive tensor power of the canonical line on the minimal
model of $Z_i$ has a meromorphic section whose divisor is
supported on the images of these interface and conductor
curves.

When this minimal model $M$ is of class VII, it has the
following restrictions. If $b_2(M)=0$, it is a Hopf surface.
If $b_2(M)>0$, there are an integer $n>0$, a nonzero
effective integral divisor $A$ and a complex-flat line
bundle $L$ such that
\[
                 K_M^{\otimes n}\otimes\OO_M(A)
                         \simeq L.                       \tag{C.CA.2}
\]
Consequently $M$ contains a global spherical shell by the
proved numerically anticanonical theorem. In particular,
the unrestricted normalization problem does not require
the classification of class VII surfaces without a
numerically anticanonical divisor. These are restrictions
on the components of the actual map $f$; no occurrence of
every surface with (C.CA.2) is asserted.
\end{statement}
\begin{proof}
After shrinking the disc, $f_*K_U$ is a locally free
rank-one sheaf. Coherence follows from properness, and
torsion-freeness follows because a holomorphic canonical
section killed by a power of the base coordinate vanishes
on a dense open subset. Its rank is one since every nearby
torus has a one-dimensional space of canonical forms.
Choose a frame and evaluate it on $U$. Over $\Delta^*$,
Grauert base change identifies its restriction with a
nonzero torus canonical form, which is nowhere zero.
The evaluation is therefore a nonzero holomorphic section
of $K_U$ with divisor $\sum_j a_jS_j$, where $a_j\ge0$.
This constructs the claimed line-bundle isomorphism.

Let $N_i=h_i^*\OO_U(S_i)$. The canonical sections defining
$S_j$ do not vanish identically on $S_i$ for $j\ne i$.
Their pullbacks define the actual effective divisors
$E_{ij}$, including any exceptional orders. Pullback of
the principal full fibre gives
\[
 N_i^{\otimes m_i}\otimes
 \OO_{Z_i}\!\left(\sum_{j\ne i}m_jE_{ij}\right)
                         \simeq\OO_{Z_i}.               \tag{C.CA.3}
\]
Adjunction for the individual reduced hypersurface $S_i$,
followed by the conductor comparison of A.5.3, gives
\[
 \begin{split}
 K_{Z_i}\otimes\OO_{Z_i}(D_i)
   &\simeq h_i^*\bigl((K_U\otimes\OO_U(S_i))|_{S_i}\bigr)\\
   &\simeq N_i^{\otimes(a_i+1)}\otimes
      \OO_{Z_i}\!\left(\sum_{j\ne i}a_jE_{ij}\right).
 \end{split}                                             \tag{C.CA.4}
\]
Raise (C.CA.4) to the $m_i$-th power and substitute
(C.CA.3). This proves (C.CA.1), with its signs and every
full-fibre multiplicity retained.

Let $\tau:Z_i\to M$ be the sequence of point blowdowns
to a minimal model. Identity (C.CA.1) provides a
meromorphic section of $K_{Z_i}^{\otimes m_i}$ with the
displayed signed divisor, less $m_iD_i$. On the open set
where $\tau$ is an isomorphism this is a meromorphic
$m_i$-canonical section on $M$. In local coordinates at
each omitted point, the coefficient of that section is a
meromorphic function on a punctured smooth surface and
extends meromorphically. Thus the section extends across
those points. Its divisor is precisely the pushforward:
exceptional curves disappear, and every other curve
keeps its order. Equivalently, at an individual blowup
the exact canonical formula
$K_{\widehat M}=\tau^*K_M+E$ removes the corresponding
$m_iE$ from the pulled-back pluricanonical divisor.
We have therefore constructed a signed integral divisor
$B$ on $M$ and an actual isomorphism
\[
                         K_M^{\otimes m_i}\simeq\OO_M(B).
                                                               \tag{C.CA.5}
\]
Every component of $B$ is the image of an interface or
intrinsic conductor curve from (C.CA.1).

Suppose now that $M$ is of class VII. If it contained no
curve, (C.CA.5) would have $B=0$ and would trivialize a
positive power of $K_M$. This contradicts
$\kappa(M)=-\infty$. When $b_2(M)=0$, choose a nonempty
curve. A positive multiple of it is homologically trivial,
since its integral class is torsion. The classification
of such effective divisors in
\cite[Remark~2.7]{DlousskyTeleman2012} gives a Hopf
surface in this zero-Betti case. This proves the first
alternative without an assertion about unknown class VII
surfaces.

For $b_2(M)>0$, let $C_1,\ldots,C_r$ be the distinct
curves in the support of $B$, and write
\[
 c_1(K_M)_{\mathbb R}=\sum_{\ell=1}^r c_\ell[C_\ell],
 \qquad c_\ell\in\Q.                                  \tag{C.CA.6}
\]
The intersection form on $H^2(M,\mathbb R)$ is negative
definite and $K_M^2=-b_2(M)$. Moreover
$K_M\cdot C_\ell\ge0$ for every curve on the minimal
surface; these standard class VII facts are stated with
their adjunction argument in \cite[\S1, pp.~4--5]{Dloussky2006}.
For completeness, adjunction gives
$K_M\cdot C=2p_a(C)-2-C^2$. A curve of arithmetic genus
zero is a smooth rational curve; minimality excludes
self-intersection $-1$, and self-intersection zero would
give zero real class and hence $K_M\cdot C=0$, contradicting
adjunction. Every other curve has arithmetic genus at
least one and $C^2\le0$, so also has nonnegative canonical
intersection.

Let $Q=(C_\ell\cdot C_k)$ and
$v_\ell=K_M\cdot C_\ell$. Then
\[
                         Qc=v,\qquad v\ge0.             \tag{C.CA.7}
\]
The matrix $Q$ is negative semidefinite, and all its
off-diagonal entries are nonnegative. Separate it into
blocks according to the connected components of the
intersection graph. For a negative-definite block, write
$c=c_+-c_-$ with disjoint nonnegative coordinate supports.
If $c_+\ne0$, then
\[
 c_+^tQc=c_+^tQc_+-c_+^tQc_-<0,
 \qquad c_+^tv\ge0,                                    \tag{C.CA.8}
\]
contradicting (C.CA.7). Thus every coefficient in this
block is nonpositive.

We also treat every singular block, rather than assuming
linear independence of the curve classes. Choose a nonzero
kernel vector $u$. Replacing it by its coordinatewise
absolute value cannot decrease its quadratic form, because
the off-diagonal entries are nonnegative. Both quadratic
forms must be zero by negative semidefiniteness. Hence
$|u|$ is again in the kernel. A zero coordinate of $|u|$
would, in the corresponding kernel equation, force every
neighboring coordinate to be zero. Connectedness would
then force the whole vector to vanish. Thus the block has
a strictly positive kernel vector. Pairing (C.CA.7) with
this vector gives $v=0$ on the block, and then $Qc=0$.
The real divisor class $\beta=\sum c_\ell[C_\ell]$ of
that block satisfies $\beta^2=0$. The ambient negative
definite form forces $\beta=0$.

Discard the singular blocks and negate the coefficients
on the negative-definite blocks. This gives an effective
rational divisor $A_{\Q}$ on the actual curves of $M$
such that
\[
                   c_1(K_M+A_{\Q})_{\mathbb R}=0.        \tag{C.CA.9}
\]
It is nonzero, since otherwise $c_1(K_M)_{\mathbb R}=0$
would contradict $K_M^2=-b_2(M)<0$.

Choose a positive multiple clearing the denominators in
$A_{\Q}$ and then another positive multiple killing the
torsion integral Chern class in (C.CA.9). We obtain an
integer $n>0$, a nonzero effective integral divisor $A$,
and a topologically trivial holomorphic line
\[
                       L=K_M^{\otimes n}\otimes\OO_M(A).
                                                               \tag{C.CA.10}
\]
We specify why this line is complex-flat. On a class VII
surface $p_g=0$, $c_2(M)=b_2(M)$ and $K_M^2=-b_2(M)$,
so Noether's formula gives $\chi(\OO_M)=0$ and
$h^1(\OO_M)=1$. Also $\dim_{\C}H^1(M,\C)=1$.
The natural map $H^1(M,\R)\to H^1(M,\OO_M)$ is
injective: if a real closed one-form $\alpha$ has
$\alpha^{0,1}=\bar\partial(p+\mathrm i q)$, with $p,q$
real, then
$\alpha=dp+\mathrm i(\bar\partial q-\partial q)$ and
closedness gives $\partial\bar\partial q=0$. The
pluriharmonic maximum principle on the compact connected
surface makes $q$ constant, so $\alpha$ is exact.
Consequently the complex-linear map
$H^1(M,\C)\to H^1(M,\OO_M)$ is an isomorphism.
The exponential sequences for constant and holomorphic
coefficients therefore give the surjection
\[
             H^1(M,\C^*)\longrightarrow\Pic^0(M).       \tag{C.CA.11}
\]
This is the actual flat-character map, taking constant
transition functions to the same holomorphic transition
functions. It supplies a complex-flat realization of $L$.

The canonical section of $\OO_M(A)$ now becomes a
nonzero section of $K_M^{-n}\otimes L$. These are exactly
the hypotheses of \cite[Theorem~0.2]{Dloussky2006},
which proves existence of a global spherical shell.
This proves (C.CA.2) and the stated normalization
restriction. The argument also shows directly why a
positive-Betti class VII surface whose every curve class
is zero in real cohomology cannot be a minimal
normalization model here: (C.CA.6) would contradict
$K_M^2=-b_2(M)$.
\end{proof}

C.2.UN.19 subsequently excludes the positive-$b_2$ alternative
altogether by the actual cyclic base-change quotient and its
descent to a Hopf surface. The canonical-support and
numerically anticanonical comparison above is retained as an
independent necessary reduction. It does not supply an example
of a positive-$b_2$ class VII component.
